%----------------------------------------------------------------------------------------
%   USEFUL COMMANDS
%----------------------------------------------------------------------------------------

\newcommand{\dipartimento}{Dipartimento di Matematica ``Tullio Levi-Civita''}

%----------------------------------------------------------------------------------------
% 	USER DATA
%----------------------------------------------------------------------------------------

% Data di approvazione del piano da parte del tutor interno; nel formato GG Mese AAAA
% compilare inserendo al posto di GG 2 cifre per il giorno, e al posto di 
% AAAA 4 cifre per l'anno
\newcommand{\dataApprovazione}{[Data di approvazione --- GG Mese AAAA]}

% Dati dello Studente
\newcommand{\nomeStudente}{Giulio}
\newcommand{\cognomeStudente}{Bottacin}
\newcommand{\matricolaStudente}{2042340}
\newcommand{\emailStudente}{giulio.bottacin@studenti.unipd.it}
\newcommand{\telStudente}{+39 340 574 7169}

% Dati del Tutor Aziendale
\newcommand{\nomeTutorAziendale}{Stefano}
\newcommand{\cognomeTutorAziendale}{Bertolin}
\newcommand{\emailTutorAziendale}{stefano.bertolin@kireygroup.com}
\newcommand{\telTutorAziendale}{[Telefono tutor aziendale]}
\newcommand{\ruoloTutorAziendale}{[Ruolo in azienda]}

% Dati dell'Azienda
\newcommand{\ragioneSocAzienda}{Kirey S.r.l.}
\newcommand{\indirizzoAzienda}{Corso Stati Uniti, 14/B, 35127 Padova (PD)}
\newcommand{\sitoAzienda}{https://www.kireygroup.com/}
\newcommand{\emailAzienda}{stefano.bertolin@kireygroup.com}
\newcommand{\partitaIVAAzienda}{CF/P.IVA 06729880960}

% Dati del Tutor Interno (Docente)
\newcommand{\titoloTutorInterno}{Prof.}
\newcommand{\nomeTutorInterno}{Tullio}
\newcommand{\cognomeTutorInterno}{Vardanega}

\newcommand{\prospettoSettimanale}{
    \begin{itemize}
        \item \textbf{Prima Settimana {[XX ore]} --- {[Titolo fase, es. Onboarding e formazione]}}
        \begin{itemize}
            \item Incontro con le persone coinvolte nel progetto;
            \item Verifica credenziali e strumenti di lavoro assegnati;
            \item Presa visione dell’infrastruttura esistente;
            \item {[Altra attività prevista]};
        \end{itemize}
        \item \textbf{Seconda Settimana {[XX ore]} --- {[Titolo fase]}}
        \begin{itemize}
            \item {[Attività prevista]};
            \item {[Attività prevista]};
        \end{itemize}
        \item \textbf{Terza Settimana {[XX ore]} --- {[Titolo fase]}}
        \begin{itemize}
            \item {[Attività prevista]};
            \item {[Attività prevista]};
        \end{itemize}
        \item \textbf{Quarta Settimana {[XX ore]} --- {[Titolo fase]}}
        \begin{itemize}
            \item {[Attività prevista]};
            \item {[Attività prevista]};
        \end{itemize}
        \item \textbf{Quinta Settimana {[XX ore]} --- {[Titolo fase]}}
        \begin{itemize}
            \item {[Attività prevista]};
            \item {[Attività prevista]};
        \end{itemize}
        \item \textbf{Sesta Settimana {[XX ore]} --- {[Titolo fase]}}
        \begin{itemize}
            \item {[Attività prevista]};
            \item {[Attività prevista]};
        \end{itemize}
        \item \textbf{Settima Settimana {[XX ore]} --- {[Titolo fase]}}
        \begin{itemize}
            \item {[Attività prevista]};
            \item {[Attività prevista]};
        \end{itemize}
        \item \textbf{Ottava Settimana {[XX ore]} --- {[Titolo fase, es. Collaudo e documentazione finale]}}
        \begin{itemize}
            \item {[Attività prevista]};
            \item {[Attività prevista]};
        \end{itemize}
    \end{itemize}
}

% Indicare il totale complessivo (deve essere compreso tra le 300 e le 320 ore)
\newcommand{\totaleOre}{[Totale ore --- tra 300 e 320]}

\newcommand{\obiettiviObbligatori}{
	 \item \underline{\textit{O01}}: [Primo obiettivo obbligatorio];
	 \item \underline{\textit{O02}}: [Secondo obiettivo obbligatorio];
	 \item \underline{\textit{O03}}: [Terzo obiettivo obbligatorio];
}

\newcommand{\obiettiviDesiderabili}{
	 \item \underline{\textit{D01}}: [Primo obiettivo desiderabile];
	 \item \underline{\textit{D02}}: [Secondo obiettivo desiderabile];
}

\newcommand{\obiettiviFacoltativi}{
	 \item \underline{\textit{F01}}: [Primo obiettivo facoltativo];
	 \item \underline{\textit{F02}}: [Secondo obiettivo facoltativo];
}